\section*{Bab \ref{ch:b3:measure} --- Teori Ukuran}

\begin{solution}{exo:b3:measure:1}
(a) Pengambilan komplemen menukar kedua kasus pendefinisinya. Gabungan
terbilang himpunan terbilang tetap terbilang; dan jika salah satu anggotanya
berkomplemen terbilang, maka gabungannya pun demikian: jadi kestabilannya berlaku. Ia
memuat himpunan tunggal, dan setiap aljabar-$\sigma$ yang memuatnya
memuat semua himpunan terbilang (lewat gabungan terbilang) beserta
komplemennya: jadi ia $\sigma(\{\text{himpunan tunggal}\})$.

(b) Tulis $\mathcal B = \sigma(\text{himpunan terbuka})$. Setiap himpunan bagian
$\R$ yang terbuka merupakan gabungan terbilang selang terbuka
berujung rasional (di sekitar setiap titik rasional \omterm{def:b3:topology:topology}{himpunan terbukanya},
ambil selang berjari-jari rasional di dalamnya), sehingga \omterm{def:b3:topology:topology}{himpunan terbukanya} $\in
\sigma(\text{selang terbuka}) \subseteq \sigma(\text{data
rasional})$. Peralihannya: $\intoo ab = \bigcup_n\intcc{a +
\frac1n}{b - \frac1n}$; $\intcc ab = \bigcap_n \intoo{a -
\frac1n}{b + \frac1n}$; $\intoc{-\infty}a = \bigcap_n
\intoo{-\infty}{a + \frac1n}$ dan sebaliknya $\intoo ab =
\intoo{-\infty}b \setminus \intoc{-\infty}a$; sedangkan untuk sinar rasional:
$\intoc{-\infty}a =
\bigcap_{q \in \Q,\, q > a}\intoc{-\infty}q$. Jadi setiap keluarganya
membangkitkan yang lain lewat operasi terbilang: keempatnya membangkitkan
$\mathcal B$.

(c) Dengan ujung yang berhingga saja, keluarganya bahkan bukan
aljabar: sebab komplemen $\intoc01$ memuat sinar tak terbatas.
Bila ujung tak berhingga diperbolehkan ($\intoc{-\infty}b$, $\intoo
a{+\infty}$) ia menjadi aljabar (sebab komplemen dan gabungan
berhingga atas gabungan semacam itu tetap demikian), tetapi bukan aljabar-$\sigma$:
sebab $\{0\} = \bigcap_n\intoc{-\frac1n}0$ bukan gabungan berhingga
selang tak degenerat.
\end{solution}

\begin{solution}{exo:b3:measure:2}
(a) $A \cup B = A \sqcup (B \setminus (A\cap B))$, sehingga
$\mu(A\cup B) = \mu(A) + \mu(B) - \mu(A\cap B)$ (dan keberhinggaannya
mengizinkan pengurangan itu). Untuk tiga himpunan: terapkan rumus dua himpunannya
dua kali,
\[
\mu(A\cup B\cup C) = \sum\mu(A) - \sum\mu(A\cap B) +
\mu(A\cap B\cap C)
\]
(dengan jumlah atas himpunan indeks yang sudah jelas).

(b) Untuk \omterm{def:b3:measure:lebesgueouter}{ukuran Lebesgue}, $A_n = [n, +\infty) \downarrow
\varnothing$, padahal $\lambda(A_n) = \infty \not\to 0$.

(c) Sebuah titik terletak di selang berpanjang $\varepsilon$:
sehingga $\lambda(\{x\}) = 0$; lalu subaditivitas terbilangnya mematikan himpunan
terbilang. Karena itu $\lambda(\Q \cap \intcc01) = 0$ dan, lewat keaditifannya,
$\lambda(\intcc01\setminus\Q) = 1$: jadi bilangan irasionallah yang membawa seluruh
panjangnya.
\end{solution}

\begin{solution}{exo:b3:measure:3}
Sinar $\intoc{-\infty}t$ membentuk sistem-$\pi$ (sebab
irisan dua di antaranya adalah yang lebih kecil) yang membangkitkan $\mathcal B(\R)$
(\cref{exo:b3:measure:1}). Himpunan $P_k = \intoc{-\infty}k$
naik menuju $\R$ dengan $\mu(P_k) \leq 1 < \infty$: sehingga
\cref{thm:b3:measure:uniqueness} berlaku, dan $\mu = \nu$ pada
$\mathcal B(\R)$. Dengan demikian fungsi distribusinya menentukan
\omterm{def:b3:measure:measure}{ukurannya}.
\end{solution}

\begin{solution}{exo:b3:measure:4}
Untuk setiap $N$ berlaku $\limsup A_n \subseteq \bigcup_{n \geq N}A_n$,
sehingga $\mu(\limsup A_n) \leq \sum_{n\geq N}\mu(A_n)$, yakni ekor
sebuah deret konvergen: lalu biarkan $N \to \infty$. Terapannya: dengan
$A_n = \{x \in \intcc01 : \abs{x - r_n} \leq 4^{-n}\}$
(dengan $r_n$ bilangan rasional ke-$n$), $\lambda(A_n) \leq 2\cdot4^{-n}$
dapat dijumlahkan: sehingga $\lambda(\limsup A_n) = 0$, yakni hampir setiap $x$
termasuk hanya berhingga banyak $A_n$. (Padahal setiap $x$ merupakan limit
bilangan rasional: intinya adalah \emph{kelajuan} $4^{-n}$-nya.)
\end{solution}

\begin{solution}{exo:b3:measure:5}
Himpunan $K = \bigcap K_n$ merupakan irisan gabungan berhingga selang
tertutup: jadi \omterm{def:b3:topology:compact}{kompak}. Pada tahap $n$ tersisa $2^n$ selang
berpanjang sama $\ell_n \leq 2^{-n}$ (sebab setiap tahapnya membelah dan
menyusutkan); dan selang $I \subseteq K$ akan terletak di dalam satu
selang tahap-$n$ untuk setiap $n$, sehingga memaksa $\ell(I) = 0$: jadi \omterm{def:b3:topology:interior}{interiornya}
kosong. \omterm{def:b3:measure:measure}{Ukuran} yang dibuang adalah $\sum_{n\geq1}2^{n-1}\cdot
4^{-n} = \frac12\sum_{n\geq1}2^{-n} = \frac12$, dengan semua pembuangannya
berupa selang terbuka yang saling lepas: jadi $\lambda(K) = \frac12$.

Ragamnya: membuang selang tengah berpanjang $\varepsilon
4^{-n}$ meninggalkan \omterm{def:b3:topology:compact}{kompak} tak padat di mana pun $K^{(\varepsilon)}$ yang
berukuran $1 - \frac\varepsilon2$. Lalu $\bigcup_m
K^{(1/m)}$ bersifat \omterm{rem:b3:complete:meagre}{kurus} (sebagai gabungan terbilang himpunan tak padat di mana pun)
dan berukuran $\geq \sup_m(1 - \frac1{2m}) = 1$: jadi \omterm{rem:b3:complete:meagre}{himpunan kurus} yang
berukuran penuh --- sedangkan komplemennya di $\intoo01$ merupakan $G_\delta$
padat berukuran $0$ (gemuk secara \omterm{def:b3:topology:topology}{topologi}, nol secara
metrik). Komplemen $K^{(\varepsilon)}$ di $\intoo01$ bersifat
terbuka, padat, dan berukuran $\frac\varepsilon2 < \varepsilon$.
\end{solution}

\begin{solution}{exo:b3:measure:6}
Dengan memotong $\intoc01$ menjadi $2^n$ translasi $\intoc0{2^{-n}}$:
$c = 2^n\,\mu(\intoc0{2^{-n}})$, sehingga $\mu(\intoc0{2^{-n}}) =
c\,2^{-n} = c\,\lambda(\intoc0{2^{-n}})$. Lewat keinvarianan translasi
dan keaditifannya, $\mu = c\lambda$ pada setiap selang
$\intoc ab$ yang $b - a$-nya rasional diadik dan $a$-nya apa pun; sedangkan
$\intoc ab$ yang umum merupakan gabungan naik selang semacam itu
($b_k \uparrow b$ lewat langkah diadik dari $a$), dan \omterm{def:b3:topology:continuity}{kekontinuan} dari
bawah memperluas kesamaannya. Selang $\intoc ab$ membentuk
sistem-$\pi$ yang membangkitkan $\mathcal B(\R)$, dengan $\intoc{-k}k
\uparrow \R$ yang berukuran berhingga (sebab $\mu(\intoc{-k}k) = 2kc$):
sehingga \cref{thm:b3:measure:uniqueness} memberikan $\mu = c\lambda$ pada
$\mathcal B(\R)$.
\end{solution}

\begin{solution}{exo:b3:measure:7}
Menurut keteraturannya (\cref{thm:b3:measure:regularity}) pilihlah $K
\subseteq A \subseteq U$ dengan $K$ \omterm{def:b3:topology:compact}{kompak}, $U$ terbuka, dan
$\lambda(U\setminus K) < \varepsilon$ (dengan kedua penghampirannya
dalam $\varepsilon/2$, sebab $\lambda(U \setminus K) =
\lambda(U\setminus A) + \lambda(A \setminus K)$). Tulis $U$ sebagai
gabungan terbilang selang terbuka $(I_n)$ yang saling lepas (yakni
komponen \omterm{def:b3:topology:topology}{himpunan terbukanya}); lalu $K$ yang \omterm{def:b3:topology:compact}{kompak} terselimuti oleh
berhingga di antaranya, $K \subseteq B = I_1\cup\dots\cup I_N \subseteq
U$. Maka $A \setminus B \subseteq A\setminus K \subseteq
U\setminus K$ dan $B \setminus A \subseteq U \setminus A
\subseteq U\setminus K$: sehingga
$\lambda(A\,\triangle\,B) \leq 2\lambda(U\setminus K)$ ---
jadi mulailah dari $\varepsilon/2$ agar mendarat di bawah $\varepsilon$.
\end{solution}

\begin{solution}{exo:b3:measure:8}
Dengan mengganti $A$ oleh $A \cap [-M, M]$ yang berukuran positif (dan suatu $M$
memang berhasil, menurut \omterm{def:b3:topology:continuity}{kekontinuan} dari bawah), andaikan $0 < \lambda(A) <
\infty$. Ambil $U \supseteq A$ yang terbuka dengan $\lambda(U) <
\frac43\lambda(A)$ lalu uraikan $U = \bigsqcup_nI_n$ atas
selang terbuka yang saling lepas: maka $\lambda(A) = \sum_n\lambda(A\cap
I_n)$. Seandainya setiap $n$ mempunyai $\lambda(A\cap I_n) \leq
\frac34\ell(I_n)$, menjumlahkannya akan memberikan $\lambda(A) \leq
\frac34\lambda(U) < \lambda(A)$: jadi ada selang $I$ yang memenuhi
$\lambda(A\cap I) > \frac34\ell(I)$. Tetapkan $B = A \cap I$ lalu misalkan
$\abs t < \frac12\ell(I)$: maka $B$ dan $B + t$ keduanya terletak di
selang $I \cup (I + t)$ yang berpanjang $< \frac32\ell(I)$. Seandainya
keduanya saling lepas: $\lambda(B) + \lambda(B + t) = 2\lambda(B)
> \frac32\ell(I)$ akan melampaui \omterm{def:b3:measure:measure}{ukuran} selang pemuatnya
--- mustahil. Jadi $B \cap (B + t) \neq \varnothing$:
sehingga suatu $x \in B$ ditulis $x = y + t$ dengan $y \in B$, dan $t = x -
y \in A - A$.
Karena itu $\intoo{-\frac{\ell(I)}2}{\frac{\ell(I)}2} \subseteq A -
A$.
\end{solution}

\begin{solution}{exo:b3:measure:9}
Translasi Vitali $(V + q)_{q\in\Q}$ mempartisi $\R$ (sebab setiap
bilangan real ekuivalen dengan tepat satu wakilnya). Andaikan semua
himpunan $B_q = A \cap (V + q)$ terukur. Dua unsur
$V + q$ mana pun berselisih bilangan irasional atau nol (sebab dua
wakil yang berbeda tak ekuivalen), sehingga $B_q - B_q$
memotong $\Q$ hanya di $\{0\}$: jadi ia tak memuat selang, dan
Steinhaus (\cref{exo:b3:measure:8}) memaksa $\lambda(B_q) = 0$.
Lalu $\lambda(A) \leq \sum_{q}\lambda(B_q) = 0$, dan itu bertentangan dengan
$\lambda(A) > 0$. Jadi ada $B_q \subseteq A$ yang tak terukur.
\end{solution}

\begin{solution}{exo:b3:measure:10}
\emph{Terukur $\Rightarrow$ penghampiran-$\varepsilon$}: lewat
keteraturan luarnya (\cref{thm:b3:measure:regularity}) ambillah
$U \supseteq A$ yang terbuka dengan $\lambda(U) \leq \lambda(A) +
\varepsilon$; lalu keterukurannya mengizinkan pengurangan
$\lambda(U\setminus A) = \lambda(U) - \lambda(A) \leq
\varepsilon$. \emph{Versi-$\varepsilon$ $\Rightarrow$
versi-$G_\delta$}: ambillah $U_n$ dengan $\lambda^*(U_n\setminus
A) < \frac1n$ lalu $G = \bigcap U_n$: sebuah $G_\delta$ dengan
$\lambda^*(G\setminus A) \leq \lambda^*(U_n\setminus A) \to 0$.
\emph{Versi-$G_\delta$ $\Rightarrow$ terukur}: sebab $G\setminus
A$ berukuran-$\lambda^*$ nol, sehingga terukur menurut \omterm{def:b3:complete:complete}{kelengkapannya}
(\cref{thm:b3:measure:caratheodory}), dan $A = G \setminus
(G\setminus A)$ terukur (sebab $G$ bersifat Borel). Jadi himpunan Lebesgue
tepat berupa ``Borel modulo himpunan nol''.
\end{solution}

\begin{solution}{exo:b3:measure:11}
(a) Himpunan $\liminf A_n = \bigcup_N\bigcap_{n\geq N}A_n$ merupakan
gabungan naik himpunan $B_N = \bigcap_{n\geq N}A_n$,
sehingga $\mu(\liminf A_n) = \lim_N\mu(B_N)$ (lewat \omterm{def:b3:topology:continuity}{kekontinuan} dari
bawah); dan $\mu(B_N) \leq \inf_{n \geq N}\mu(A_n)$, yang
limitnya $\liminf\mu(A_n)$. Untuk $\limsup$-nya: terapkan hal
yang sama pada komplemennya di dalam lingkungan berukuran berhingga
$U = \bigcup A_n$ --- sebab \omterm{def:b3:topology:continuity}{kekontinuan} dari atas pada
$C_N = \bigcup_{n\geq N}A_n$ yang turun menuntut $\mu(C_1)
\leq \mu(U) < \infty$, dan memberikan $\mu(\limsup A_n) =
\lim\mu(C_N) \geq \limsup\mu(A_n)$.

(b) Selang berjalan $A_n = \intcc n{n+1}$: setiap titiknya
termasuk paling banyak dua $A_n$ dan akhirnya tak termasuk satu pun,
sehingga $\limsup A_n = \varnothing$; padahal $\mu(A_n) = 1$. Jadi
$\limsup\mu(A_n) = 1 > 0 = \mu(\limsup A_n)$: tanpa selubung
berukuran berhingga, ketaksamaan kedua pada (a) gagal
--- sebab massanya lolos ke tak hingga, tempat yang tak dapat
ditangkap himpunan tetap mana pun.

(c) Jika $\sum\mu(A_n) < \infty$: maka $\mu(C_N) \leq
\sum_{n\geq N}\mu(A_n) \to 0$ dan $\limsup A_n =
\bigcap C_N$ berukuran $\leq \inf\mu(C_N) = 0$. Untuk kasus
monotonnya: yang naik adalah \omterm{def:b3:topology:continuity}{kekontinuan} dari bawah; sedangkan yang turun dengan
$\mu(A_1) < \infty$ adalah \omterm{def:b3:topology:continuity}{kekontinuan} dari atas ---
keduanya sudah dibuktikan pada sifat dasar \cref{ch:b3:measure}; dan
contoh penyangkal $A_n = \intco n\infty$ (yang turun ke
$\varnothing$ dengan $\mu \equiv \infty$) menunjukkan bahwa keberhinggaannya
kembali menjadi hal yang mendasar.
\end{solution}

\begin{solution}{exo:b3:measure:12}
(a) Himpunan $E_{k,N}$ naik terhadap $N$ (sebab kendalanya
berkurang), dan setiap $x$ akhirnya memenuhi $\abs{f_n(x)
- f(x)} \leq \frac1k$ untuk setiap $n \geq N(x)$ (menurut kekonvergenan
titik demi titiknya): jadi $\bigcup_NE_{k,N} = X$. Lalu lewat \omterm{def:b3:topology:continuity}{kekontinuan} dari bawah:
$\mu(E_{k,N}) \to \mu(X) < \infty$, sehingga $\mu(X\setminus
E_{k,N}) \to 0$; lalu pilihlah $N_k$ yang sesuai.

(b) Misalkan $A = \bigcap_kE_{k,N_k}$: maka $\mu(X\setminus A) \leq
\sum_k\varepsilon2^{-k} = \varepsilon$. Pada $A$: untuk setiap
$k$, setiap $n \geq N_k$ memenuhi $\sup_A\abs{f_n - f} \leq
\frac1k$ --- dan itu tepat kekonvergenan seragam pada $A$.

(c) Untuk gundukan berjalannya, kekonvergenan seragam pada $A$ menuntut
$A$ akhirnya menghindari setiap $\intcc n{n+1}$ --- lebih
tepatnya $\sup_A\abs{f_n} < \frac12$ memaksa $A \cap \intcc
n{n+1}$ kosong untuk $n$ besar, sehingga $X \setminus A$
memuat sebuah ekor $\bigcup_{n\geq n_0}\intcc n{n+1}$ yang
berukuran tak berhingga. Pada (a), keberhinggaannya mengubah
``$E_{k,N}\nearrow X$'' menjadi ``\omterm{def:b3:measure:measure}{ukuran} komplemennya
menuju $0$'': sebab \omterm{def:b3:topology:continuity}{kekontinuan} dari atas memerlukan awal yang berhingga,
dan pada ruang berukuran tak berhingga pelolosan ke tak hingga persis
merupakan hal yang tak dapat dilihatnya.
\end{solution}

\begin{solution}{pb:b3:measure:1}
\textbf{1.} Lewat induksi. \omterm{def:b3:topology:continuity}{Kekontinuannya}: ketiga rumusnya bersesuaian di
sambungannya (sebab $\frac12c_n(1) = \frac12$ dan $\frac12 +
\frac12c_n(0) = \frac12$); dan setiap kepingannya \omterm{def:b3:topology:continuity}{kontinu}.
Kemonotonan dan nilai batasnya diwarisi. Untuk
taksiran kontraksinya: pada $[0,\frac13]$, $\abs{c_{n+1} - c_n}(x)
= \frac12\abs{c_n - c_{n-1}}(3x) \leq \frac12\norm{c_n -
c_{n-1}}_\infty$; pada sepertiga tengahnya selisihnya $0$; dan pada
sepertiga kanannya, sama seperti kirinya.

\textbf{2.} Berlaku $\norm{c_{n+1} - c_n}_\infty \leq
2^{-n}\norm{c_1 - c_0}_\infty$: sehingga deret pertambahannya
konvergen seragam, jadi $c_n \to c$ secara seragam; lalu $c$ bersifat
\omterm{def:b3:topology:continuity}{kontinu}, tak turun, dengan $c(0) = 0$, $c(1) = 1$ (semuanya
terpelihara oleh limit seragam), dan mengambil limit pada
rekursi pendefinisinya menunjukkan bahwa $c$ sendiri memenuhi ketiga
identitas serupa-dirinya.

\textbf{3.} Menurut identitas tengahnya, $c \equiv \frac12$ pada
$\intcc{\frac13}{\frac23}$, yakni celah pertamanya. Setiap celah $C$ merupakan
peta celah pertama itu di bawah komposisi kedua kontraksi
afin $x \mapsto \frac x3$ dan $x\mapsto\frac{x +
2}3$; dan identitasnya mengangkut kekonstanannya sejalan dengan itu (dengan
nilainya berupa bilangan rasional diadik). Untuk rumus angkanya, ambillah $x =
\sum_n 2b_n3^{-n} \in C$: jika $b_1 = 0$ maka $x \in
[0,\frac13]$ dan $c(x) = \frac12c(3x)$ dengan $3x$ berangka
$(b_2, b_3, \dots)$; sedangkan jika $b_1 = 1$ maka $x \in [\frac23, 1]$
dan $c(x) = \frac12 + \frac12c(3x - 2)$, dengan pergeseran yang sama. Lewat
induksi, $N$ angka biner pertama $c(x)$ adalah $b_1,
\dots, b_N$ untuk setiap $N$: jadi $c(x) = \sum_nb_n2^{-n}$.

\textbf{4.} Di luar $C$, $c$ konstan secara lokal: sehingga dapat diturunkan
dengan turunan $0$. Dan karena $\lambda(C) = 0$
(\cref{ex:b3:measure:cantor}), maka $c' = 0$ hampir di mana-mana. Namun
$c(1) - c(0) = 1$: sebab teorema fundamental dalam bentuk $\mathcal
C^1$-nya menuntut $c$ dapat diturunkan \emph{di mana-mana}
dengan turunan yang \omterm{def:b3:topology:continuity}{kontinu} (atau setidaknya terintegralkan, ditambah \omterm{def:b3:topology:continuity}{kekontinuan}
mutlak --- lihat \cref{ch:b3:lebesgue}); padahal $c$ tak
dapat diturunkan di titik $C$, dan yang lebih mendasar, $c$
menggagalkan \emph{\omterm{def:b3:topology:continuity}{kekontinuan} mutlak}: sebab ia mendaki pada himpunan nol.

\textbf{5.} Diberikan $y = \sum_n\beta_n2^{-n} \in \intcc01$
(dengan $\beta_n \in \{0,1\}$), titik $x = \sum_n2\beta_n3^{-n}
\in C$ mempunyai $c(x) = y$ menurut pertanyaan 3: jadi $c(C) = \intcc01$, yakni himpunan
berukuran $1$ --- sehingga himpunan nol $C$ membawa, lewat $c$,
seluruh selangnya.

\textbf{6.} Fungsi $h$ \omterm{def:b3:topology:continuity}{kontinu} dan naik sejati (sebab $x$ demikian
dan $c$ tak turun); lalu $h(0) = 0$, $h(1) = 1$, sehingga menurut
teorema nilai antara $h$ merupakan bijeksi \omterm{def:b3:topology:continuity}{kontinu} pada
$\intcc01$; dan bijeksi \omterm{def:b3:topology:continuity}{kontinu} dari \omterm{def:b3:topology:compact}{ruang kompak} ke \omterm{def:b3:topology:hausdorff}{ruang
Hausdorff} merupakan \omterm{def:b3:topology:continuity}{homeomorfisma}
(\cref{cor:b3:topology:compacthomeo}).

\textbf{7.} Pada sebuah celah $(u, v)$ (berpanjang $\ell$), $c$ konstan,
sehingga $h$ bersifat afin berkemiringan $\frac12$: jadi $h((u,v))$ merupakan
selang berpanjang $\frac\ell2$. Celahnya saling lepas dan $h$
injektif: sehingga petanya saling lepas, dengan \omterm{def:b3:measure:measure}{ukuran} total
$\frac12\sum\ell = \frac12(1 - \lambda(C)) = \frac12$.

\textbf{8.} Berlaku $h(\intcc01) = \intcc01$ dan $h(C)$ bersifat \omterm{def:b3:topology:compact}{kompak}
(sebagai peta \omterm{def:b3:topology:continuity}{kontinu}), sehingga terukur, dengan
\[
\lambda\bigl(h(C)\bigr) = 1 -
\lambda\bigl(h(\intcc01\setminus C)\bigr) = 1 - \tfrac12 =
\tfrac12 .
\]
Jadi sebuah \omterm{def:b3:topology:continuity}{homeomorfisma} dapat menggembungkan himpunan nol menjadi berukuran $\frac12$:
sehingga ``\omterm{def:b3:measure:measure}{ukuran} \omterm{def:b3:topology:topology}{topologi}'' (kategori, dimensi) dan ``\omterm{def:b3:measure:measure}{ukuran}'' diangkut
\omterm{def:b3:topology:continuity}{homeomorfisma} dengan cara yang sangat berbeda --- hanya yang
pertama yang merupakan invarian \omterm{def:b3:topology:topology}{topologi}.

\textbf{9.} Himpunan $Z = h^{-1}(W) \subseteq h^{-1}(h(C)) = C$ mempunyai
$\lambda^*(Z) \leq \lambda(C) = 0$: jadi himpunan nol, sehingga
terukur Lebesgue menurut \omterm{def:b3:complete:complete}{kelengkapannya}
(\cref{thm:b3:measure:caratheodory}).

\textbf{10.} Misalkan $\varphi$ \omterm{def:b3:topology:continuity}{kontinu} dan $\mathcal D = \{B
: \varphi^{-1}(B) \in \mathcal B\}$. Prapeta komutatif dengan
komplemen dan gabungan terbilang, sehingga $\mathcal D$ merupakan
aljabar-$\sigma$; dan ia memuat \omterm{def:b3:topology:topology}{himpunan terbukanya} (menurut \omterm{def:b3:topology:continuity}{kekontinuannya}):
jadi $\mathcal D \supseteq \mathcal B$ --- sehingga prapeta himpunan Borel
di bawah \omterm{def:b3:topology:continuity}{pemetaan kontinu} bersifat Borel.

\textbf{11.} Seandainya $Z$ bersifat Borel, terapkan pertanyaan 10 pada
$\varphi = h^{-1}$ yang \omterm{def:b3:topology:continuity}{kontinu}: maka $\varphi^{-1}(Z) = h(Z) = W$
akan bersifat Borel, sehingga terukur Lebesgue --- dan itu bertentangan dengan
pemilihan $W$. Jadi $Z \in \mathcal L \setminus \mathcal B$:
sehingga aljabar-$\sigma$ Lebesgue memuat milik Borel secara sejati.

\textbf{12.} Fungsi $g = \mathbf 1_Z$ terukur Lebesgue (sebab $Z \in
\mathcal L$) dan $\varphi = h^{-1}$ \omterm{def:b3:topology:continuity}{kontinu}, tetapi
$(g\circ\varphi)^{-1}(\{1\}) = \varphi^{-1}(Z) = W$ tidak
terukur: jadi $g \circ \varphi$ tak terukur Lebesgue.
Kehati-hatian yang diperlukan: ``fungsi terukur Lebesgue'' berarti prapeta
himpunan \emph{Borel} mendarat di $\mathcal L$; sedangkan mengomposisikannya menuntut
prapeta himpunan \emph{Lebesgue} tetap Lebesgue, dan itu tak
diberikan \omterm{def:b3:topology:continuity}{kekontinuan} (di sini $\varphi^{-1}(Z) \notin
\mathcal L$ walaupun $\varphi$ sebuah \omterm{def:b3:topology:continuity}{homeomorfisma}).

\textbf{13.} Rantainya, beserta saksi kesejatiannya:
\[
\{\text{terbilang}\} \subsetneq \{\text{nol Borel}\}
\subsetneq \{\text{nol Lebesgue}\} \subsetneq \mathcal L
\subsetneq \mathcal P(\R),
\qquad
\{\text{nol Borel}\} \subsetneq \mathcal B \subsetneq
\mathcal L .
\]
Saksinya: $C$ bersifat Borel, nol, dan tak terbilang (untuk jurang pertama); $Z$
bersifat nol Lebesgue tetapi bukan Borel (untuk jurang kedua dan, di dalam $\mathcal B
\subsetneq \mathcal L$, jurang keenam); himpunan Cantor gemuk bersifat Borel,
tak padat di mana pun, dan berukuran positif (yang memisahkan himpunan nol dari
himpunan Borel); dan himpunan Vitali $V$ tak berada di $\mathcal L$ (untuk jurang terakhir).
Jadi \omterm{def:b3:measure:measure}{ukuran}, \omterm{def:b3:topology:topology}{topologi}, dan kardinalitas mengiris $\mathcal P(\R)$ menurut
garis yang sungguh berbeda.

\textbf{14.} Untuk \omterm{def:b3:measure:outer}{ukuran luarnya}: $\mu_F^*(\varnothing) = 0$ (selimuti
dengan selang yang melenyap), kemonotonannya jelas, dan subaditivitas
terbilangnya menyusul lewat penyambungan selimut yang optimal dalam
$\varepsilon2^{-k}$, persis seperti untuk $\lambda^*$. Selimut satu selangnya
memberikan $\mu_F^*(\intoc ab) \leq F(b) - F(a)$. Sebaliknya misalkan
$\intoc ab \subseteq \bigcup_k\intoc{a_k}{b_k}$. Menurut \omterm{def:b3:topology:continuity}{kekontinuan}
$F$, ambillah $b_k' > b_k$ dengan $F(b_k') \leq F(b_k) +
\varepsilon2^{-k}$, dan $a' \in \intoo ab$ dengan $F(a') \leq
F(a) + \varepsilon$. \omterm{def:b3:topology:compact}{Kompak} $\intcc{a'}b$ terselimuti oleh
$\intoo{a_k}{b_k'}$ yang terbuka: sehingga berhingga di antaranya sudah cukup, dan
argumen perantaian pada \cref{thm:b3:measure:lebesgue} (yakni berjalan dari
$a'$ ke $b$ melalui selang yang bertindih, sambil meneleskopkan
pertambahan $F$, dengan kemonotonannya menyerap tindihannya) melahirkan $F(b)
- F(a') \leq \sum_k(F(b_k') - F(a_k)) \leq \sum_k(F(b_k) -
F(a_k)) + \varepsilon$. Lalu biarkan $\varepsilon \to 0$: diperoleh kesamaan.
(Dan \omterm{def:b3:topology:continuity}{kekontinuan} $F$-lah yang mengizinkan pembukaan selangnya dengan
biaya-$F$ sekecil apa pun.)

\textbf{15.} Cukup dibuktikan bahwa setiap setengah garis $H_t =
\intoc{-\infty}t$ terukur menurut Carathéodory, sebab
himpunan terukurnya membentuk aljabar-$\sigma$
(\cref{thm:b3:measure:caratheodory}) dan setengah garis membangkitkan
$\mathcal B$. Diberikan $A$ dan sebuah selimut optimal dalam $\varepsilon$,
sebut $(\intoc{a_k}{b_k})$, atas $A$: maka setiap selangnya terpecah sebagai
$\intoc{a_k}{t\wedge b_k} \cup \intoc{t \vee a_k}{b_k}$ (dengan satu
kepingannya mungkin kosong), dan biaya-$F$-nya berjumlah tepat
$F(b_k) - F(a_k)$; lalu kepingan pertamanya menyelimuti $A \cap H_t$, dan
kepingan keduanya menyelimuti $A \setminus H_t$. Karena itu $\mu_F^*(A\cap H_t) +
\mu_F^*(A\setminus H_t) \leq \mu_F^*(A) + \varepsilon$, sedangkan
ketaksamaan sebaliknya adalah subaditivitasnya. Dengan membatasi
\omterm{def:b3:measure:measure}{ukuran} yang dihasilkan pada $\mathcal B$: diperolehlah \omterm{def:b3:measure:measure}{ukuran}
Lebesgue--Stieltjes $\mu_F$.

\textbf{16.} Berlaku $\mu(\R) = \lim_n(F(n) - F(-n)) = 1 - 0 = 1$
(lewat \omterm{def:b3:topology:continuity}{kekontinuan} \omterm{def:b3:measure:measure}{ukurannya} sepanjang $\intoc{-n}n$). Pada sebuah celah
$\intoo uv$ dari $C$, $c$ konstan, sehingga setiap subselang setengah
terbukanya berukuran-$\mu$ nol dan demikian pula celahnya (lewat gabungan
terbilang); dan di luar $\intcc01$, $c$ juga konstan. Karena itu
$\mu(\R\setminus C) = 0$ dan $\mu(C) = 1$, sedangkan $\lambda(C) =
0$ (\cref{ex:b3:measure:cantor}): jadi masing-masing dari $\mu, \lambda$
dibawa oleh himpunan yang oleh yang lain dinyatakan nol --- yakni saling
singular.

\textbf{17.} Berlaku $\mu(\{x\}) = \lim_{\delta\downarrow0}
\mu(\intoc{x-\delta}x) = \lim(c(x) - c(x-\delta)) = 0$ menurut
\omterm{def:b3:topology:continuity}{kekontinuan} $c$: jadi tak beratom. Sehingga $\mu$ merupakan \omterm{def:b3:measure:measure}{ukuran}
peluang tanpa atom yang dibawa sebuah \omterm{def:b3:topology:compact}{kompak} berukuran nol Lebesgue ---
bukan menyebar-berkepadatan seperti pembatasan $\lambda$,
bukan pula beratom seperti \omterm{def:b3:measure:measure}{ukuran} cacah: yakni spesies ketiga.

\textbf{18.} Kepingan $C_\varepsilon$ merentang sebuah selang
terner $I_\varepsilon$ berpanjang $3^{-m}$, dan Bagian I
pertanyaan 3 menunjukkan bahwa sepanjang $I_\varepsilon$ tangganya
mendaki tepat $2^{-m}$ (sebab $m$ angka biner pertama $c$
dibekukan pada $\varepsilon$, sedangkan sisanya menyapu segalanya).
Karena itu $\mu(C_\varepsilon) = \mu(I_\varepsilon) =
c(\text{ujung kanan}) - c(\text{ujung kiri}) = 2^{-m}$: jadi
silinder angka berkedalaman $m$ semuanya bermassa $2^{-m}$, yakni hukum $m$ koin
adil.

\textbf{19.} Ruas kanannya mendefinisikan \omterm{def:b3:measure:measure}{ukuran} Borel
$\nu = \frac12\,(x \mapsto \tfrac x3)_*\mu + \frac12\,(x
\mapsto \tfrac{x+2}3)_*\mu$ yang dinilaikan di $A$ --- yakni \omterm{def:b3:measure:measure}{ukuran}
peluang. Secara titik demi titik, dengan $c$ yang diperluas secara global, mudah diperiksa
kasus demi kasus ($x \leq 0$; ketiga sepertiganya; $x \geq 1$)
identitas
\[
c(x) = \tfrac12\,c(3x) + \tfrac12\,c(3x - 2),
\]
misalnya pada $\intcc{1/3}{2/3}$: $\frac12\cdot1 + \frac12\cdot0 =
\frac12 = c(x)$. Karena itu menilaikan $\nu$ pada $\intoc ab$
memberikan $c(b) - c(a) = \mu(\intoc ab)$, sedangkan dua \omterm{def:b3:measure:measure}{ukuran} berhingga
yang bersesuaian pada sistem-$\pi$ berisi selang setengah terbuka berimpit
pada $\mathcal B$ (\cref{thm:b3:measure:uniqueness}): jadi $\nu =
\mu$.

\textbf{20.} Lewat induksi pada $n$: $c_0(1-x) = 1 - c_0(x)$,
dan jika $c_n$ mempunyai kesimetrian itu, maka untuk $x \in
\intcc0{1/3}$: $c_{n+1}(1 - x) = \frac12 + \frac12c_n(3(1-x)
- 2) = \frac12 + \frac12c_n(1 - 3x) = \frac12 + \frac12(1 -
c_n(3x)) = 1 - c_{n+1}(x)$; sedangkan sepertiga tengahnya bercermin di sekitar
$\frac12$; dan sepertiga kanannya merupakan kasus kirinya yang dicerminkan. Pada
limitnya, $c(1-x) = 1 - c(x)$. Untuk peta majunya: $(s_*\mu)(\intoc ab)
= \mu(\intco{1-b}{1-a}) = c(1-a) - c(1-b)$ (sebab $\mu$ tak beratom,
pertanyaan 17, sehingga kesepakatan batasnya tak berbiaya) $= (1 -
c(a)) - (1 - c(b)) = \mu(\intoc ab)$: jadi $s_*\mu = \mu$ menurut
ketunggalannya.

\textbf{21.} Misalkan $X \sim \mu$ (integral fungsi \omterm{def:b3:topology:continuity}{kontinu}
terhadap $\mu$ ada sebagai limit jumlah bertipe Riemann
atas kepingan berkedalaman $m$, yang masing-masing bermassa $2^{-m}$, dengan
galat pencuplikan $\leq \operatorname{osc} \leq
\norm{f'}_\infty3^{-m}$; dan \cref{ch:b3:lebesgue} akan
menyistematiskannya). Kesimetriannya: $1 - X \sim X$, sehingga $\E X =
\frac12$. Keserupaan dirinya: $X$ berhukum $\frac Y3$ dengan
peluang $\frac12$ dan berhukum $\frac{Y+2}3$ dengan peluang
$\frac12$, dengan $Y \sim \mu$, sehingga
\[
\E X^2 = \frac12\,\frac{\E Y^2}9 + \frac12\,
\frac{\E Y^2 + 4\E Y + 4}{9} = \frac{2\E X^2 + 6}{18},
\]
sehingga $\E X^2 = \frac38$ dan $\operatorname{Var}X = \frac38
- \frac14 = \frac18$. Sedangkan hukum seragam beragam
$\frac1{12} < \frac18$: jadi massa Cantor memeluk ujungnya.

\textbf{22.} Himpunan $C$ tertutup dan $\mu(C) = 1$. Jika sebuah $I$ yang terbuka
memotong $C$ di $x$, maka kepingan berkedalaman $m$ yang memuat $x$ menyusut
ke $x$, sehingga ada $I_\varepsilon \subseteq I$ dan $\mu(I) \geq
2^{-m} > 0$: jadi tak ada himpunan tertutup yang lebih kecil yang dapat membawa $\mu$. Sedangkan titik
di luar $C$ berpersekitaran celah yang berukuran-$\mu$ sama dengan $0$. Karena itu
$\operatorname{supp}\mu = C$ tepat.

\textbf{23.} Satu gejala, dua dialek. Bagian I mengatakan bahwa
pertumbuhan $c$ tak terlihat oleh turunannya: $c' = 0$ hampir di mana-mana,
dengan seluruh pendakiannya memusat pada himpunan nol $C$. Bagian V
mengatakan bahwa \omterm{def:b3:measure:measure}{ukuran} terkaitnya $\mu_c$ menaruh seluruh massanya pada
himpunan nol yang sama: $\mu_c \perp \lambda$, sehingga tak ada kepadatan $f \geq
0$ yang dapat memenuhi $\mu_c(A) = \int_Af\,\dd\lambda$ --- sebab kepadatan
memaksa kelenyapan pada himpunan nol-$\lambda$. Kamusnya: bagi yang
tak turun dan terbatas, $F \leftrightarrow$ \omterm{def:b3:measure:measure}{ukuran} berhingga
$\mu_F$ (pertanyaan 14--15); $F$ yang merupakan integral turunannya
$\leftrightarrow$ $\mu_F$ yang berkepadatan (yakni kasus ``\omterm{def:b3:topology:continuity}{kontinu}
mutlak''); dan secara umum $F'$, yang ada hampir di mana-mana
bagi $F$ yang monoton (lewat teorema pendiferensialan Lebesgue, yang melampaui
bab ini), hanya memulihkan bagian kepadatannya. Tangganya
merupakan ujung ekstremnya: \omterm{def:b3:topology:continuity}{kontinu}, dengan turunan $0$ hampir di mana-mana ---
sehingga \omterm{def:b3:measure:measure}{ukurannya} \emph{singular murni}, dan teorema
fundamental kalkulus, jauh dari gagal secara kebetulan, gagal persis
sebesar massa singularnya $\mu_c(\R) = 1$.

\textbf{24.} Tetapkan $x < y$ di $\intcc01$ lalu pilihlah $m \geq 0$
dengan $3^{-(m+1)} < y - x \leq 3^{-m}$. Sepanjang sembarang selang
triadik $\intcc{k3^{-m}}{(k+1)3^{-m}}$ tangganya mendaki
paling banyak $2^{-m}$: sebab selang semacam itu entah salah satu dari $2^m$
kepingan $C_m$, yang pendakiannya tepat $2^{-m}$ (pertanyaan
18), entah termuat di \omterm{def:b3:topology:interior}{penutup} satu celah pada suatu
tahap $\leq m$, tempat $c$ konstan. Dan karena $y - x \leq
3^{-m}$, selang $\intcc xy$ memotong paling banyak dua
selang triadik berkedalaman $m$ yang berurutan, sehingga
\[
c(y) - c(x) \leq 2\cdot2^{-m}
= 2\bigl(3^{-m}\bigr)^{s}
< 2\bigl(3(y - x)\bigr)^{s}
= 4\,(y - x)^{s},
\]
dengan memakai $3^{s} = 2$ dan $3^{-m} < 3(y - x)$. Keoptimalannya:
ujung $u < v$ sebuah kepingan Cantor berkedalaman $m$ memenuhi $v - u =
3^{-m}$ dan $c(v) - c(u) = 2^{-m} = (v - u)^{s}$; sehingga batas Hölder
$\abs{c(v) - c(u)} \leq K(v - u)^{t}$ dengan $t > s$ akan
memaksa $K \geq 2^{-m}3^{mt} = (3^{t}/2)^{m} \to \infty$ (sebab
$3^{t} > 3^{s} = 2$), dan setiap subselang $\intcc01$
memuat kepingan semacam itu, sehingga kegagalannya bersifat lokal pula.
Bentuk \omterm{def:b3:measure:measure}{ukurannya}: $\mu\bigl(\intcc{x-r}{x+r}\bigr) \leq
c(x + r) - c(x - r) \leq 4(2r)^{s} = 4\cdot2^{s}r^{s} \leq
8r^{s}$ (dengan nilai $c$ yang diperluas ke $\R$ seperti pada pertanyaan 16;
sebab $\mu$ tak beratom, pertanyaan 17).

\textbf{25.} Tulis $S_0(x) = \frac x3$ dan $S_1(x) =
\frac{x+2}3$; maka hipotesisnya mengatakan $\nu = \frac12(S_0)_*\nu +
\frac12(S_1)_*\nu$. Dengan mengiterasikannya $m$ kali,
\[
\nu = 2^{-m}\sum_{w \in \{0,1\}^m}(S_w)_*\nu,
\qquad S_w = S_{w_1}\circ\dots\circ S_{w_m}.
\]
Karena $\nu$ dibawa oleh $\intcc01$ dan $S_w(\intcc01) =
I_w$, yakni kepingan Cantor berkedalaman $m$ yang diindeks $w$, maka setiap
$(S_w)_*\nu$ merupakan \omterm{def:b3:measure:measure}{ukuran} peluang yang dibawa $I_w$; sedangkan
$2^m$ kepingannya berupa selang tertutup saling lepas berpasangan yang berpanjang
$3^{-m}$. Tetapkan $\intoc ab$ lalu misalkan $N_m$ banyaknya
kepingan $I_w \subseteq \intoc ab$. Sebuah kepingan yang memotong $\intoc ab$
tanpa termuat di dalamnya harus memuat $a$ atau $b$, sedangkan sebuah
titik terletak paling banyak di satu kepingan, sehingga
\[
2^{-m}N_m \leq \nu(\intoc ab) \leq 2^{-m}N_m + 2\cdot2^{-m}.
\]
Ketaksamaan rangkap yang sama berlaku bagi $\mu$ (sebab pertanyaan 19 memberikan
iterasi yang identik), dengan cacahan $N_m$ yang sama. Karena itu
$\abs{\nu(\intoc ab) - \mu(\intoc ab)} \leq 2^{-m+1} \to 0$:
sehingga $\nu$ dan $\mu$ bersesuaian pada sistem-$\pi$ berisi selang setengah
terbuka, dan keduanya \omterm{def:b3:measure:measure}{ukuran} peluang, sehingga
\cref{thm:b3:measure:uniqueness} memberikan $\nu = \mu$. Jadi
\omterm{def:b3:measure:measure}{ukuran} tangganya merupakan \emph{satu-satunya} titik tetap
skema perata-rataan dua pemetaan itu --- yakni pernyataan setingkat \omterm{def:b3:measure:measure}{ukuran} atas
keserupaan diri $C$.

\end{solution}
